\documentclass[11pt,a4paper]{article}
\usepackage[top=1.6cm,bottom=1.4cm,left=1.3cm,right=1.3cm,headsep=8pt,footskip=16pt]{geometry}
\usepackage{enumitem}
\usepackage{amssymb}
\usepackage{array}
\usepackage{xcolor}
\usepackage{fancyhdr}
\usepackage{multicol}
\usepackage{polyglossia}
\setdefaultlanguage[calendar=gregorian, numerals=maghrib]{arabic}
\setotherlanguage{english}
\IfFontExistsTF{Amiri}{%
  \newfontfamily\arabicfont[Script=Arabic, Scale=1.08]{Amiri}
}{%
  \newfontfamily\arabicfont[Script=Arabic, Scale=1.08, Path=../fonts/, Extension=.ttf,
    UprightFont=*-Regular, BoldFont=*-Bold, ItalicFont=*-Italic,
    BoldItalicFont=*-BoldItalic]{Amiri}
}
\IfFontExistsTF{AlBattar}{%
  \newfontfamily\khatt[Script=Arabic, Scale=1.15]{AlBattar}
}{%
  \newfontfamily\khatt[Script=Arabic, Scale=1.15, Path=../fonts/,
    BoldFont=ae_AlBattar.ttf]{ae_AlBattar.ttf}
}
\renewcommand{\headrulewidth}{0.8pt}
\renewcommand{\footrulewidth}{0.8pt}
\fancyhf{}
\rhead{\khatt المدرسة العليا للأساتذة ببوزريعة}
\lhead{\khatt ملحقة بومرداس}
\lfoot{\khatt الأستاذ: د. بن عبد الرزاق يوسف}
\cfoot{\thepage/2}
\pagestyle{fancy}
\parindent 0pt
\setlength{\columnsep}{0.6cm}
\setlength{\columnseprule}{0.4pt}
\setlist[enumerate,1]{label=\textbf{\arabic*.}, leftmargin=*, itemsep=5pt, topsep=2pt}
\begin{document}

\begin{center}
{\Large\khatt\textbf{الاختبار الثامن في الذكاء الاصطناعي}}\\[2pt]
{\normalsize أسئلة متعددة الاختيارات ــ الفصول من 1 إلى 6 ــ 15 سؤالًا}
\end{center}
\vspace{-4pt}\rule{\textwidth}{0.6pt}\\
{\small\textbf{التعليمة:} لكل سؤال إجابة صحيحة واحدة فقط من بين خمسة اختيارات.
اكتب \textbf{حرف} الإجابة الصحيحة في جدول الإجابات أسفل الصفحة الثانية.
لكل سؤال نقطة واحدة (المجموع 15 نقطة).}
\vspace{-2pt}
\begin{multicols}{2}

\noindent\begin{minipage}{\linewidth}
\hangindent=0.7cm\hangafter=1\makebox[0.7cm][r]{\textbf{1.}}\textbf{البيانات الجيدة يجب أن تكون:}
\begin{itemize}[leftmargin=0.95cm, labelsep=4pt, itemsep=0pt, topsep=2pt, parsep=0pt]
  \item[\textbf{أ)}] مكررة كثيرًا
  \item[\textbf{ب)}] قليلة ومن فئة واحدة
  \item[\textbf{ج)}] قديمة ومتشابهة
  \item[\textbf{د)}] متنوعة وممثِّلة
  \item[\textbf{هـ)}] من مصدر واحد
\end{itemize}
\end{minipage}\par\vspace{7pt plus 4pt}
\noindent\begin{minipage}{\linewidth}
\hangindent=0.7cm\hangafter=1\makebox[0.7cm][r]{\textbf{2.}}\textbf{استعمال الأداة لفهم قاعدة نحوية صعبة استعمال:}
\begin{itemize}[leftmargin=0.95cm, labelsep=4pt, itemsep=0pt, topsep=2pt, parsep=0pt]
  \item[\textbf{أ)}] مرفوض
  \item[\textbf{ب)}] ممنوع دائمًا
  \item[\textbf{ج)}] انتحال
  \item[\textbf{د)}] غير مسؤول
  \item[\textbf{هـ)}] مقبول
\end{itemize}
\end{minipage}\par\vspace{7pt plus 4pt}
\noindent\begin{minipage}{\linewidth}
\hangindent=0.7cm\hangafter=1\makebox[0.7cm][r]{\textbf{3.}}\textbf{ما المخطط المفاهيمي؟}
\begin{itemize}[leftmargin=0.95cm, labelsep=4pt, itemsep=0pt, topsep=2pt, parsep=0pt]
  \item[\textbf{أ)}] صورة فوتوغرافية للقسم وللتلاميذ
  \item[\textbf{ب)}] رسم يربط المفاهيم الأساسية بأسهم
  \item[\textbf{ج)}] قائمة مرتبة لمراجع البحث ومصادره
  \item[\textbf{د)}] جدول يجمع علامات التلاميذ في القسم
  \item[\textbf{هـ)}] تسجيل صوتي لشرح الأستاذ للدرس
\end{itemize}
\end{minipage}\par\vspace{7pt plus 4pt}
\noindent\begin{minipage}{\linewidth}
\hangindent=0.7cm\hangafter=1\makebox[0.7cm][r]{\textbf{4.}}\textbf{لماذا الطلب «اشرح النعت لتلميذ في السنة الأولى متوسط، بلغة بسيطة، مع ثلاثة أمثلة» أفضل من «اشرح النعت»؟}
\begin{itemize}[leftmargin=0.95cm, labelsep=4pt, itemsep=0pt, topsep=2pt, parsep=0pt]
  \item[\textbf{أ)}] لأنه يمنع الأداة من الهلوسة نهائيًا
  \item[\textbf{ب)}] لأنه يحدد المستوى واللغة وعدد الأمثلة
  \item[\textbf{ج)}] لأنه يجعل الأداة تنسخ الكتاب المدرسي
  \item[\textbf{د)}] لأنه يُلزم الأداة بالإجابة بالإنجليزية
  \item[\textbf{هـ)}] لأنه أطول، والطول وحده يضمن الجودة
\end{itemize}
\end{minipage}\par\vspace{7pt plus 4pt}
\noindent\begin{minipage}{\linewidth}
\hangindent=0.7cm\hangafter=1\makebox[0.7cm][r]{\textbf{5.}}\textbf{لدينا 500 صورة هواتف و500 صورة حواسيب و100 صورة ألواح رقمية. ماذا نتوقع؟}
\begin{itemize}[leftmargin=0.95cm, labelsep=4pt, itemsep=0pt, topsep=2pt, parsep=0pt]
  \item[\textbf{أ)}] يتعرف على الألواح أفضل لأن صورها قليلة
  \item[\textbf{ب)}] لا يخطئ النموذج أبدًا مع أي فئة منها
  \item[\textbf{ج)}] يخطئ النموذج مع الهواتف أكثر من غيرها
  \item[\textbf{د)}] يخطئ النموذج كثيرًا مع الألواح الرقمية
  \item[\textbf{هـ)}] يخطئ النموذج بالتساوي مع الفئات الثلاث
\end{itemize}
\end{minipage}\par\vspace{7pt plus 4pt}
\noindent\begin{minipage}{\linewidth}
\hangindent=0.7cm\hangafter=1\makebox[0.7cm][r]{\textbf{6.}}\textbf{أي مما يلي ليس طريقة لتقليل الانحياز؟}
\begin{itemize}[leftmargin=0.95cm, labelsep=4pt, itemsep=0pt, topsep=2pt, parsep=0pt]
  \item[\textbf{أ)}] استعمال بيانات من فئة واحدة لتسريع التدريب
  \item[\textbf{ب)}] مراجعة أخطاء النموذج بعد اختباره وتحليلها
  \item[\textbf{ج)}] جمع بيانات متنوعة تمثل جميع الفئات والحالات
  \item[\textbf{د)}] اختبار النموذج على حالات مختلفة ومتنوعة
  \item[\textbf{هـ)}] الموازنة بين عدد الأمثلة في كل فئة من الفئات
\end{itemize}
\end{minipage}\par\vspace{7pt plus 4pt}
\noindent\begin{minipage}{\linewidth}
\hangindent=0.7cm\hangafter=1\makebox[0.7cm][r]{\textbf{7.}}\textbf{أي عنصر يجيب عن السؤال: «ما الشروط التي يجب احترامها؟»}
\begin{itemize}[leftmargin=0.95cm, labelsep=4pt, itemsep=0pt, topsep=2pt, parsep=0pt]
  \item[\textbf{أ)}] المهمة
  \item[\textbf{ب)}] الدور
  \item[\textbf{ج)}] السياق
  \item[\textbf{د)}] التنسيق
  \item[\textbf{هـ)}] القيود
\end{itemize}
\end{minipage}\par\vspace{7pt plus 4pt}
\noindent\begin{minipage}{\linewidth}
\hangindent=0.7cm\hangafter=1\makebox[0.7cm][r]{\textbf{8.}}\textbf{ما الانتحال العلمي؟}
\begin{itemize}[leftmargin=0.95cm, labelsep=4pt, itemsep=0pt, topsep=2pt, parsep=0pt]
  \item[\textbf{أ)}] التحقق من المعلومة في مصدر موثوق قبل نقلها
  \item[\textbf{ب)}] التصريح باستعمال أداة ذكاء اصطناعي في البحث
  \item[\textbf{ج)}] ذكر المصادر بدقة في آخر البحث وفي الهوامش
  \item[\textbf{د)}] نسبة عمل أو أفكار الغير إلى النفس دون ذكر المصدر
  \item[\textbf{هـ)}] تلخيص نص بأسلوب شخصي مع ذكر مصدره الأصلي
\end{itemize}
\end{minipage}\par\vspace{7pt plus 4pt}
\noindent\begin{minipage}{\linewidth}
\hangindent=0.7cm\hangafter=1\makebox[0.7cm][r]{\textbf{9.}}\textbf{معيار «الجهة» في تقييم المصدر يجيب عن السؤال:}
\begin{itemize}[leftmargin=0.95cm, labelsep=4pt, itemsep=0pt, topsep=2pt, parsep=0pt]
  \item[\textbf{أ)}] هل هدفه الإعلام أم البيع؟
  \item[\textbf{ب)}] من كتبه؟ هل هو مختص في المجال؟
  \item[\textbf{ج)}] من نشره؟ جامعة أم موقع مجهول؟
  \item[\textbf{د)}] متى نُشر؟ هل ما زال صالحًا؟
  \item[\textbf{هـ)}] هل يذكر مراجعه ومصادره بدقة؟
\end{itemize}
\end{minipage}\par\vspace{7pt plus 4pt}
\noindent\begin{minipage}{\linewidth}
\hangindent=0.7cm\hangafter=1\makebox[0.7cm][r]{\textbf{10.}}\textbf{رفع أستاذ صورة جدول وطلب شرحًا مكتوبًا له. ما المُدخل وما المُخرَج؟}
\begin{itemize}[leftmargin=0.95cm, labelsep=4pt, itemsep=0pt, topsep=2pt, parsep=0pt]
  \item[\textbf{أ)}] المدخل نص، والمخرج صوت
  \item[\textbf{ب)}] المدخل صورة، والمخرج نص
  \item[\textbf{ج)}] المدخل نص، والمخرج صورة
  \item[\textbf{د)}] المدخل فيديو، والمخرج صوت
  \item[\textbf{هـ)}] المدخل صوت، والمخرج نص
\end{itemize}
\end{minipage}\par\vspace{7pt plus 4pt}
\noindent\begin{minipage}{\linewidth}
\hangindent=0.7cm\hangafter=1\makebox[0.7cm][r]{\textbf{11.}}\textbf{أي العبارات الآتية صحيحة؟}
\begin{itemize}[leftmargin=0.95cm, labelsep=4pt, itemsep=0pt, topsep=2pt, parsep=0pt]
  \item[\textbf{أ)}] كل ما يقدّمه الذكاء الاصطناعي صحيح إذا كان الطلب دقيقًا
  \item[\textbf{ب)}] الذكاء الاصطناعي التوليدي لا ينتج إلا نصوصًا مكتوبة
  \item[\textbf{ج)}] قد تتغير الإجابة حسب صياغة الطلب، فنتحقق من المهم
  \item[\textbf{د)}] الأدوات المشهورة لا تقبل الطلبات المكتوبة بالعربية
  \item[\textbf{هـ)}] الطلب الأقصر والأبسط يعطي دائمًا الإجابة الأفضل
\end{itemize}
\end{minipage}\par\vspace{7pt plus 4pt}
\noindent\begin{minipage}{\linewidth}
\hangindent=0.7cm\hangafter=1\makebox[0.7cm][r]{\textbf{12.}}\textbf{الأمر: «أنشئ خمسة أسئلة فهم حول النص الآتي، مع الإجابات». أي عنصرين غائبين عنه؟}
\begin{itemize}[leftmargin=0.95cm, labelsep=4pt, itemsep=0pt, topsep=2pt, parsep=0pt]
  \item[\textbf{أ)}] الدور والسياق
  \item[\textbf{ب)}] المهمة والتنسيق
  \item[\textbf{ج)}] السياق والمهمة
  \item[\textbf{د)}] المهمة والقيود
  \item[\textbf{هـ)}] القيود والتنسيق
\end{itemize}
\end{minipage}\par\vspace{7pt plus 4pt}
\noindent\begin{minipage}{\linewidth}
\hangindent=0.7cm\hangafter=1\makebox[0.7cm][r]{\textbf{13.}}\textbf{ما الفرق بين الإفصاح والشفافية؟}
\begin{itemize}[leftmargin=0.95cm, labelsep=4pt, itemsep=0pt, topsep=2pt, parsep=0pt]
  \item[\textbf{أ)}] لا فرق بينهما إطلاقًا، فهما اسمان لمبدأ واحد
  \item[\textbf{ب)}] الإفصاح يخص الأداة، والشفافية تخص الشركة المطوِّرة
  \item[\textbf{ج)}] الشفافية إخفاء الاستعمال، والإفصاح كشفه للقارئ
  \item[\textbf{د)}] الشفافية ذكر المصادر الورقية، والإفصاح ذكر المواقع
  \item[\textbf{هـ)}] الإفصاح تصريح بالاستعمال، والشفافية توضيح كيفيته
\end{itemize}
\end{minipage}\par\vspace{7pt plus 4pt}
\noindent\begin{minipage}{\linewidth}
\hangindent=0.7cm\hangafter=1\makebox[0.7cm][r]{\textbf{14.}}\textbf{طلب طالب من أداة قائمة مراجع حول موضوعه، فأعطته قائمة منسقة. كيف يتعامل معها؟}
\begin{itemize}[leftmargin=0.95cm, labelsep=4pt, itemsep=0pt, topsep=2pt, parsep=0pt]
  \item[\textbf{أ)}] يتحقق من وجود كل مرجع، لأن الأداة قد تهلوس
  \item[\textbf{ب)}] يحذف المراجع العربية ويحتفظ بالأجنبية فقط
  \item[\textbf{ج)}] يستشهد بها مباشرة لأنها منسقة تنسيقًا جيدًا
  \item[\textbf{د)}] يعتمد المرجع الأول فقط ويحذف باقي القائمة
  \item[\textbf{هـ)}] يعدّها كلها مصادر أكاديمية محكّمة وموثوقة
\end{itemize}
\end{minipage}\par\vspace{7pt plus 4pt}
\noindent\begin{minipage}{\linewidth}
\hangindent=0.7cm\hangafter=1\makebox[0.7cm][r]{\textbf{15.}}\textbf{اقترحت الأداة مخططًا جيدًا لكنه أغفل «شبه الجملة» من أنواع الخبر. أي سؤال مراجعة يكشف ذلك؟}
\begin{itemize}[leftmargin=0.95cm, labelsep=4pt, itemsep=0pt, topsep=2pt, parsep=0pt]
  \item[\textbf{أ)}] هل الخط المستعمل في المخطط جميل وواضح؟
  \item[\textbf{ب)}] هل الصورة الأصلية للدرس واضحة بما يكفي؟
  \item[\textbf{ج)}] هل يحتوي المخطط على كل الأفكار الأساسية؟
  \item[\textbf{د)}] هل أضافت الأداة مفهومًا غير موجود في النص؟
  \item[\textbf{هـ)}] هل المخطط طويل بما يكفي ليملأ الصفحة؟
\end{itemize}
\end{minipage}\par\vspace{7pt plus 4pt}
\end{multicols}
\vfill
\begin{center}
{\khatt\textbf{جدول الإجابات}}\\[3pt]
\renewcommand{\arraystretch}{1.6}
\begin{tabular}{|>{\centering\arraybackslash}m{1.9cm}|*{15}{>{\centering\arraybackslash}m{0.62cm}|}}
\hline
\textbf{رقم السؤال} & \textbf{1} & \textbf{2} & \textbf{3} & \textbf{4} & \textbf{5} & \textbf{6} & \textbf{7} & \textbf{8} & \textbf{9} & \textbf{10} & \textbf{11} & \textbf{12} & \textbf{13} & \textbf{14} & \textbf{15} \\ \hline
\textbf{الإجابة} &  &  &  &  &  &  &  &  &  &  &  &  &  &  &  \\[8pt] \hline
\end{tabular}

\vspace{8pt}
\renewcommand{\arraystretch}{2.2}
\begin{tabular}{|>{\centering\arraybackslash}m{7.5cm}|>{\centering\arraybackslash}m{3.2cm}|>{\centering\arraybackslash}m{3.2cm}|}
\hline
\textbf{الاسم واللقب} & \textbf{الفوج} & \textbf{العلامة} \\ \hline
 & & \hfill /15 \\ \hline
\end{tabular}
\end{center}
\end{document}
